% Includable section. Requires amsmath, amssymb, amsthm, hyperref.
% Bibliography entries are supplied in lower_bounds.bib.
\section{Unconditional lower bounds}
\label{sec:lower-bounds}

The lower bounds in this section do not assume the correctness of any proposed
Diophantine definition.  In fact, they hold over every number field.
Let $K$ be a number field and let $\mathcal O_K$ be its ring of integers.
For a polynomial
\[
 F(T,Y_1,\ldots,Y_r)\in K[T,Y_1,\ldots,Y_r]
\]
write
\[
 S_F=\{t\in K:\text{there exist }y_1,\ldots,y_r\in K
                    \text{ with }F(t,y_1,\ldots,y_r)=0\}.
\]
The variable $T$ is the input; the $r$ variables $Y_i$ are the witnesses.
The degree is the total degree, including $T$.

\begin{theorem}\label{thm:lower-bounds}
If $S_F=\mathcal O_K$, then
\[
                    r\geq 2\qquad\text{and}\qquad\deg F\geq 4.
\]
In particular, a Diophantine definition of $\mathbb Z$ in $\mathbb Q$
by one polynomial needs at least three total variables and has degree at least
four.  Thus a valid quartic construction has the smallest possible degree,
regardless of its number of variables.
\end{theorem}

We prove the two assertions separately.  The variable bound uses the integral
form of Hilbert irreducibility.  The cubic obstruction uses Koll\'ar's theorem
on unirational cubic hypersurfaces; the treatment of cones and reducible
polynomials is included explicitly.

\subsection{One witness and Hilbert irreducibility}

We use the following standard form of Hilbert irreducibility.  Given finitely
many irreducible polynomials $f_i(T,Y)\in K(T)[Y]$, each of positive degree in
$Y$, and finitely many forbidden values of $T$, there are infinitely many
$a\in\mathcal O_K$ outside those values for which all the specializations are
defined, preserve their $Y$-degrees, and remain irreducible in $K[Y]$.
See \cite[Chapter~4, Section~4]{lbSchinzel2000}; the integral number-field
form is also recalled in \cite[Introduction]{lbMuller2001}.
The restriction to integral specializations is essential here: the
Hilbertianity of $K$ alone, without that strengthening, is not the assertion
being applied.

\begin{proposition}\label{prop:one-witness}
Let $G\in K[T,Y]$.  If $S_G$ contains all but finitely many elements of
$\mathcal O_K$, then $S_G$ contains all but finitely many elements of $K$.
\end{proposition}
\begin{proof}
The conclusion is immediate if $G=0$.  If $G\neq0$ and has degree zero in
$Y$, then $S_G$ is finite, contrary to the hypothesis.  Otherwise factor in
$K(T)[Y]$:
\[
                G=c(T)\prod_{i=1}^{s}f_i(T,Y)^{e_i},
\]
where $c(T)\neq0$, the $e_i$ are positive integers, and the $f_i$ are monic
irreducible polynomials of positive $Y$-degree.
If some $f_i$ is linear, it has a root $u(T)\in K(T)$.  The identity
$G(T,u(T))=0$ then specializes for all but finitely many $t\in K$,
proving the conclusion.

If every $f_i$ has $Y$-degree at least two, integral Hilbert irreducibility
provides infinitely many $a\in\mathcal O_K$ at which $c(a)\neq0$ and all
$f_i(a,Y)$ remain irreducible of degree at least two.  None has a root in
$K$, so $G(a,Y)$ has no root in $K$.  This contradicts the hypothesis.
\end{proof}

Since $K\setminus\mathcal O_K$ is infinite, the proposition rules out one
witness.  For example, $1/n\notin\mathcal O_K$ for every integer $n\ge2$,
because a rational algebraic integer is an integer.  With no witness a
nonzero polynomial in $T$ has only finitely many roots, whereas the zero
polynomial defines all of $K$.  This proves the variable assertion of
Theorem~\ref{thm:lower-bounds}.

\subsection{A topological obstruction to rational parametrizations}

Let $r_1$ and $r_2$ denote the numbers of real embeddings and conjugate pairs
of complex embeddings of $K$.  Put
\[
 K_{\infty}=\mathbb R^{r_1}\times\mathbb C^{r_2},
\]
viewed as a real vector space.  The diagonal embedding
$\iota:K\longrightarrow K_{\infty}$ has dense image, while
$\iota(\mathcal O_K)$ is a discrete lattice.  Indeed, an integral basis of
$K$ identifies these images with the images of $\mathbb Q^{[K:\mathbb Q]}$
and $\mathbb Z^{[K:\mathbb Q]}$ under an invertible real linear map.

\begin{lemma}\label{lem:rational-function-integral}
Let $R\in K(U_1,\ldots,U_m)$ be nonconstant, and let $U$ be a nonempty
Zariski-open subset of $\mathbb A^m_K$ on which $R$ is defined.
Then $R(U(K))$ is not contained in $\mathcal O_K$.
\end{lemma}
\begin{proof}
Assume to the contrary that all these values lie in $\mathcal O_K$.
After shrinking $U$ we may take $U$ to be a principal open subset, defined by
one nonzero polynomial.  Apply every archimedean embedding of $K$ to its
coefficients and to those of $R$.  This gives a continuous map, on a nonempty
Euclidean open subset of $K_{\infty}^{m}$,
\[
 R_{\infty}:\prod_{v\mid\infty}K_v^m
       \longrightarrow\prod_{v\mid\infty}K_v=K_{\infty}.
\]
Choose a small connected open ball $B$ in its domain.  The diagonal image of
$K^m$ is dense in $B$.  By hypothesis its image under $R_{\infty}$ lies in
the closed discrete set $\iota(\mathcal O_K)$; by continuity the same is true
on all of $B$.  A continuous map from a connected set to a discrete set is
constant.  Each archimedean component of $R$ is therefore constant on an
open set, so its numerator minus that constant times its denominator is the
zero polynomial.  This forces $R$ to be constant, a contradiction.
\end{proof}

For $K=\mathbb Q$, this argument is simply the density of $\mathbb Q$ in
$\mathbb R$ and the discreteness of $\mathbb Z$.  For a general number field
one must use all archimedean embeddings together; $\mathcal O_K$ need not be
discrete in one selected embedding.

An integral variety $X/K$ is called \emph{unirational over $K$} if some
projective space admits a dominant rational map to $X$, defined over $K$.

\begin{corollary}\label{cor:unirational-integral}
If $X$ is unirational over $K$ and $f\in K(X)$ is nonconstant, then the values
of $f$ at the $K$-points where it is defined cannot all belong to
$\mathcal O_K$.
\end{corollary}
\begin{proof}
Compose $f$ with a dominant rational parametrization of $X$, and restrict
the source to a suitable nonempty open subset of affine space.  Dominance
makes the resulting rational function nonconstant.  Apply
Lemma~\ref{lem:rational-function-integral}.
\end{proof}

\subsection{Quadratic polynomials}

\begin{lemma}\label{lem:quadratic-integral}
If $\deg F\le2$ and $S_F\subseteq\mathcal O_K$, then $S_F$ has at most two
elements.
\end{lemma}
\begin{proof}
Suppose that $a,b,c\in K^{r+1}$ are zeros of $F$ with three distinct input
coordinates.  If these points are collinear, the restriction of $F$ to their
line has degree at most two and three distinct roots.  Thus the entire
line is contained in $F=0$.  The input is nonconstant on that line and takes
all values in $K$, a contradiction.

Otherwise the three points span an affine plane $L$ over $K$.  The input
coordinate is a nonconstant affine linear function on $L$.  If
$F|_L=0$, a nonvertical line in $L$ again gives a contradiction.  If
$F|_L$ factors into affine linear factors over $K$, one of its at most two
lines contains two of the selected points, hence has nonconstant input;
this gives the same contradiction.  A nonzero linear polynomial cannot
vanish at three noncollinear points.

The remaining case is an irreducible quadratic polynomial on $L$.
It defines an irreducible conic with a smooth $K$-point: a singular conic
has at most one singular point, so the three distinct points cannot all be
singular.  Projection from a smooth $K$-point parametrizes the conic by
rational functions over $K$.  The input is nonconstant on this conic,
for otherwise an irreducible quadratic curve would be contained in an
affine line.  Its composition with the parametrization is a nonconstant
rational function taking only values in $\mathcal O_K$, contrary to
Lemma~\ref{lem:rational-function-integral}.
\end{proof}

\subsection{Cubic cones and the sharp degree obstruction}

We use the following published theorem, in its characteristic-zero case.
\begin{theorem}[Koll\'ar]\label{thm:kollar-input}
Let $K$ be a perfect field and let $X\subseteq\mathbb P^{n+1}_K$ be a
geometrically integral cubic hypersurface of dimension $n\ge2$ which is not a cone.
If $X(K)\neq\varnothing$, then $X$ is unirational over $K$.
\end{theorem}
This is the geometrically integral case of Theorem~2 of
\cite{lbKollar2002}; the theorem there is stated for irreducible cubics.
We handle geometric reducibility separately below, so no convention about
\emph{irreducible} is needed.  Its cone exclusion must not be omitted: a cone over a smooth plane cubic need not be unirational.

For a homogeneous cubic $H$ on a finite-dimensional $K$-vector space $E$,
define its vertex subspace by
\[
 W=\{v\in E:D_vH=0\text{ as a polynomial on }E\},
 \qquad D_vH(x)=\left.\frac{d}{ds}H(x+sv)\right|_{s=0}.
\]
The equations defining $W$ are linear over $K$.  In characteristic zero,
$D_vH=0$ is equivalent to $H(x+sv)=H(x)$ for all $x,s$.
Choosing a basis of $W$ and extending it to a basis of $E$ therefore makes
$H$ independent of the coordinates along $W$.
Consequently $H$ descends to a homogeneous cubic $\overline H$ on $E/W$.
By maximality of $W$, this quotient cubic has no nonzero vertex direction
and is not a cone.  These statements remain true after extending scalars:
the kernel of a matrix over $K$ commutes with field extension.

\begin{proof}[Proof of the degree assertion in Theorem~\ref{thm:lower-bounds}]
Suppose $S_F=\mathcal O_K$ and $\deg F\le3$.
The polynomial $F$ is neither zero nor a nonzero constant.
If $F$ is reducible, each of its irreducible factors has degree at most two.
In that case, the projection is the union of the
projections of finitely many linear or quadratic factors.  Every such
projection is contained in $\mathcal O_K$, hence finite by
Lemma~\ref{lem:quadratic-integral}.  This cannot give $\mathcal O_K$.
The same lemma excludes $\deg F\le2$.
We may therefore suppose that $F$ is irreducible of degree three.

Homogenize $F$ to an irreducible cubic
\[
 H(X_0,X_1,\ldots,X_{r+1})
      =X_0^3 F(X_1/X_0,\ldots,X_{r+1}/X_0).
\]
Here $X_1/X_0$ is the input function.
We first exclude geometric reducibility.  Since the characteristic is zero,
$H$ has no repeated geometric components.  Irreducibility over $K$ makes
the Galois group act transitively on them.  Their degrees are equal and
sum to three; thus if $H$ is geometrically reducible, it is a product of
three conjugate linear forms.  A $K$-rational zero belongs to every one of
these hyperplanes, because it is fixed by Galois.  Consequently the entire
line joining any two $K$-rational zeros lies in all three hyperplanes.
Taking affine zeros with inputs $0$ and $1$ would produce every input in
$K$, a contradiction.  We may therefore assume $H$ is geometrically
integral.

Let $W$ be its vertex subspace.
For any $v\in W$, any affine rational zero $x$ normalized by $x_0=1$,
and any $s\in K$ with $1+s v_0\neq0$, the point $x+sv$ is a zero of $H$.
Its input is
\[
                       \frac{x_1+s v_1}{1+s v_0}.
\]
If $v_1-x_1v_0\neq0$, this is a nonconstant M\"obius function of $s$;
its image omits at most one element of $K$.  It cannot be contained in
$\mathcal O_K$, because $K\setminus\mathcal O_K$ is infinite.
Hence $v_1=x_1v_0$ at every such point.  Apply this to zeros with input
$0$ and input $1$, whose existence follows from $S_F=\mathcal O_K$.
It follows that
\[
                         v_0=v_1=0\qquad(v\in W).
\]

Thus both $X_0$ and $X_1$ descend to independent linear forms on $E/W$.
Choose coordinates on the quotient beginning with these forms.  The
irreducible quotient cubic $\overline H$ is not a cone, and its affine
chart $X_0=1$ has exactly the same input projection as $H$.
Indeed, every rational point on the quotient lifts through a $K$-linear
splitting of $E\to E/W$, preserving $X_0$ and $X_1$, and $H$ is constant
along the fibers.

If the projective quotient hypersurface has dimension at least two,
Theorem~\ref{thm:kollar-input} makes it unirational, since it has a
$K$-point with $X_0=1$.  Its input function $X_1/X_0$ is nonconstant,
as its image contains all of the infinite set $\mathcal O_K$.
Corollary~\ref{cor:unirational-integral} now gives a nonintegral input,
a contradiction.

If the quotient has dimension one, its affine chart is defined by a
polynomial in the input and just one further coordinate.  Its projection
cannot be $\mathcal O_K$ by Proposition~\ref{prop:one-witness}.
If the quotient has dimension zero, the affine chart has only finitely
many points and likewise cannot project onto $\mathcal O_K$.
These cases exhaust the possibilities, proving $\deg F\ge4$.
\end{proof}

\begin{remark}[Scope of the optimality statement]
The conclusion concerns one polynomial equation with all witnesses in the
field, and total degree in all variables.  It does not prohibit a system
of quadratic equations: a sum of their squares over $\mathbb Q$ is quartic.
It also does not determine the smallest possible number of witnesses at
degree four, or the smallest degree among definitions attaining the
unknown minimum number of witnesses.  The unconditional numerical
bounds established here are two witnesses and degree four; no example
meeting both bounds is asserted.
\end{remark}

\subsection{A stronger one-witness obstruction}

For completeness there is a useful strengthening which uses Siegel's
theorem rather than just Hilbert irreducibility.
In its rational-function form, Siegel's theorem says that if $C/K$ is a
smooth projective geometrically integral curve of positive genus and
$f\in K(C)$ is nonconstant, then
\[
              \{P\in C(K):f(P)\text{ is defined and belongs to }
                       \mathcal O_K\}
\]
is finite.  One can derive this from the usual theorem on affine integral
points by taking $C$ minus the poles of $f$: the map to $\mathbb A^1$
defined by $f$ is finite, and after inverting finitely many primes its
coordinate functions are integral over $\mathcal O_{K,S}[f]$.
This reduction is also used in
\cite[Section~2.1]{lbMuller2001}; see \cite[Chapter~8]{lbLang1983}.

\begin{proposition}\label{prop:one-witness-finite}
If $G\in K[T,Y]$ satisfies $S_G\subseteq\mathcal O_K$, then $S_G$ is finite.
Thus no infinite subset of $\mathcal O_K$ is the projection of the
$K$-zeros of a single polynomial with just one witness.
\end{proposition}
\begin{proof}
Factor $G$ and treat its finitely many irreducible factors separately.
Factors depending only on $T$ give finite projections.  For another factor,
let $C$ be its plane curve.  If $C$ is not geometrically irreducible, its
$K$-points lie in the intersections of distinct geometric components,
a finite set.  Otherwise pass to its smooth projective normalization
$\widetilde C$.  Away from the finitely many singular points of $C$,
rational points lift to $\widetilde C(K)$, and the input defines a
nonconstant rational function $t$ on $\widetilde C$.

If $\widetilde C$ has positive genus, Siegel's theorem gives only finitely
many points for which $t$ is integral.  If its genus is zero and it has
any rational point, then $\widetilde C\simeq\mathbb P^1_K$.
The input then becomes a nonconstant rational function on
$\mathbb P^1_K$, and on a nonempty open subset all its values would be
integral, contradicting Lemma~\ref{lem:rational-function-integral}.
If it has no rational point, only the finitely many previously excluded
singular points remain.  Each case gives a finite projection.
\end{proof}

\begin{corollary}\label{cor:cubic-finite}
Let $F\in K[T,Y_1,\ldots,Y_r]$ have total degree at most three.
If $S_F\subseteq\mathcal O_K$, then $S_F$ is finite.
\end{corollary}
\begin{proof}
For an infinite projection the proof of
Theorem~\ref{thm:lower-bounds} applies unchanged: choose any two distinct
input values in the geometric-reducibility argument and in the vertex
elimination.  In the curve case use
Proposition~\ref{prop:one-witness-finite} instead of
Proposition~\ref{prop:one-witness}.  The other cases already give a
contradiction to infinitude.
\end{proof}
